% models-core: evidence by model for the useful models, 02 (+09), 11 (+01), 16, 17. Figures and equations moved to picture.tex (2026-10-07).
% Numbers: scratchpad fact sheets 02_09_kinetics, 11_01_16_fields (latest scored round of each card).
\section{Evidence by model: the useful models}\label{sec:models}

\begin{figure*}[tp]
\includegraphics[width=\textwidth]{js/model_periods.pdf}
\caption{Where each physics model was tested. Rows: the primary model of each card, in the order of Secs.~\ref{sec:models}--\ref{sec:rest} (useful, partly useful, not useful; no card takes model 05, 07 or 12 as primary). Columns: goal periods by regime, split at scaffold changes (e.g.\ 36a--c), then natural experiments (NE). Each cell stacks the verdicts of that model's cards in that period; its height grows with the square root of the number of tests. Right: total tests and goal periods per model. Look at model 15, tested in regime III only, and at the NE columns: most models were run on most periods, but few on the interventions.}
\label{fig:modelperiods}
\end{figure*}

% fig_gallery: one small schematic per model family (Sec. IV). Few nodes, few words.
% Only the useful or partly useful models get a schematic (Vivian, 2026-10-07): 02, 11, 16, 17, 03. No Potts, no 01 on its own.
% Ising: vermillion/blue nodes (spin states). Vector spins: arrows colored by angle (cyclic 'twilight' map, sampled). Field: orange.
\definecolor{tw25}{HTML}{647DBC}\definecolor{tw40}{HTML}{6066B6}\definecolor{tw50}{HTML}{5F55AF}
\definecolor{tw60}{HTML}{5E43A5}\definecolor{tw70}{HTML}{5B3298}\definecolor{tw85}{HTML}{531D7A}\definecolor{tw100}{HTML}{421257}
\definecolor{pottsA}{HTML}{009E73}
\begin{figure*}[t]
\centering
\resizebox{0.86\textwidth}{!}{%
\begin{tikzpicture}[font=\sffamily\small,
  ag/.style={circle,draw=inkc,line width=0.5pt,minimum size=5mm,inner sep=0pt},
  rl/.style={-{Stealth[length=1.5mm]},draw=inkc!70,line width=0.5pt},
  ttl/.style={anchor=west,font=\sffamily\small\bfseries},
  lab/.style={font=\sffamily\footnotesize,text=inkc}]
% (a) Ising, M02/01
\node[ttl] at (-0.3,3.3) {(a) M02 kinetic Ising};
\node[ag,fill=spinup!85] (i1) at (0.3,2.3) {}; \node[ag,fill=spinup!85] (i2) at (1.4,2.6) {};
\node[ag,fill=spindn!85] (i3) at (2.4,1.9) {}; \node[ag,fill=spinup!85] (i4) at (0.6,1.0) {};
\node[ag,fill=spindn!85] (i5) at (1.8,0.7) {};
\draw[rl] (i1)--(i2); \draw[rl] (i2)--(i3); \draw[rl] (i4)--(i5); \draw[rl] (i5)--(i3);
\node[lab,align=center] at (1.3,-0.45) {flip at own call;\\links act at read-out};
% (b) vector spins, M11
\begin{scope}[xshift=3.6cm]
\node[ttl] at (-0.3,3.3) {(b) M11 vector spins};
\draw[-{Stealth[length=2.6mm]},draw=fieldc,line width=2.4pt] (2.65,0.5) -- (3.05,2.75);
\node[lab,text=fieldc,anchor=west] at (2.95,2.9) {goal};
\foreach \x/\y/\a/\c in {0.3/2.1/75/tw40, 1.3/2.25/105/tw85, 0.5/1.0/85/tw60, 1.6/1.3/65/tw25, 1.1/0.3/95/tw70, 2.1/2.0/80/tw50} {
  \fill[inkc!25] (\x,\y) circle (0.08);
  \draw[-{Stealth[length=2mm]},draw=\c,line width=1.6pt] (\x,\y) -- ++(\a:0.7);}
\node[lab,align=center] at (1.4,-0.45) {aligned to the field,\\weakly to each other};
\end{scope}
% (c) Langevin, M16
\begin{scope}[xshift=7.6cm]
\node[ttl] at (-0.3,3.3) {(c) M16 Langevin};
\draw[->,draw=nullc] (0,0.5) -- (3.0,0.5) node[lab,anchor=north east] at (3.0,0.45) {calls};
\draw[->,draw=nullc] (0,0.5) -- (0,2.9);
\draw[draw=inkc,line width=0.9pt,domain=0:2.9,smooth,samples=40,variable=\t] plot (\t,{0.5+1.9*exp(-\t/1.6)});
\draw[draw=spindn,line width=1.1pt,domain=0:2.9,smooth,samples=40,variable=\t] plot (\t,{0.5+1.9*exp(-\t/0.18)});
\node[lab,anchor=west] at (0.9,1.85) {well};
\node[lab,anchor=west,text=spindn] at (0.25,0.85) {kick};
\node[lab,align=center] at (1.5,-0.45) {two rates,\\ratio $\approx$15};
\end{scope}
% (d) collective modes, M17
\begin{scope}[xshift=11.3cm]
\node[ttl] at (-0.3,3.3) {(d) M17 modes};
\draw[->,draw=nullc] (0,0.5) -- (3.1,0.5) node[lab,anchor=north east] at (3.1,0.45) {eigenvalue};
\fill[nullc!45,domain=0.15:1.45,smooth,variable=\x] (0.15,0.5) -- plot (\x,{0.5+2.2*sqrt(max(0,(\x-0.15)*(1.45-\x)))}) -- (1.45,0.5) -- cycle;
\draw[dashed,draw=inkc] (1.55,0.5) -- (1.55,2.6) node[lab,anchor=south] {edge};
\fill[spindn] (2.6,0.5) rectangle (2.72,1.6);
\node[lab,anchor=south,text=spindn] at (2.66,1.6) {mode};
\node[lab,align=center] at (1.5,-0.45) {noise bulk;\\one mode above};
\end{scope}
% (e) branching, M03
\begin{scope}[xshift=15.0cm]
\node[ttl] at (-0.3,3.3) {(e) M03 branching};
\node[ag,fill=pottsA!80,minimum size=4mm] (r) at (1.4,2.5) {};
\node[ag,fill=pottsA!50,minimum size=4mm] (c1) at (0.8,1.5) {};
\node[ag,fill=pottsA!50,minimum size=4mm] (c2) at (2.0,1.5) {};
\node[ag,fill=pottsA!30,minimum size=4mm] (g1) at (2.0,0.6) {};
\draw[rl] (r)--(c1); \draw[rl] (r)--(c2); \draw[rl] (c2)--(g1);
\node[lab,align=center] at (1.4,-0.45) {$\hat R\approx0.2$:\\trees die out};
\end{scope}
\end{tikzpicture}}
\caption{The useful model families, from the most to the least useful. (a)~Kinetic Ising with read-out coupling (M02): each agent is a two-state spin (two colors) that updates only at its own calls; links act at the read-out call. (b)~Vector spins (M11): each agent's content is an arrow (color = angle); the goal is a strong field (orange), so the arrows align with it and only weakly with each other. (c)~Langevin relaxation (M16): a read kick decays about 15 times faster than the agent's own well. (d)~Collective modes (M17): a mode counts only above the calibrated noise edge. (e)~Branching (M03): idea trees die out, $\hat R\approx0.2$.}
\label{fig:gallery}
\end{figure*}


Section~\ref{sec:picture} stated the four useful models and their best evidence. This section gives the full evidence for each, and the next two sections cover the other models (Fig.~\ref{fig:gallery}). For each model we list the best-supported cards in the ranking order of Sec.~\ref{sec:faith}, what failed, and a verdict. Credences are the current judged values (Sec.~\ref{sec:faith}).


\textbf{Where each model was tested.} Coverage is not uniform (Fig.~\ref{fig:modelperiods}). Each card runs its replication layer on every goal period with enough data, so most models were tested on 31--41 of the 49 goal-period columns: model 02 has 575 tests over 41 periods, model 11 has 438 over 35, and model 01 has 419 over 40. Model 15 is the exception: its seven cards test regime III only (10 periods, 89\% of its tests). Regime I has the most periods and so the most tests. But coverage is not evidence. The headline numbers come from narrower windows, which the \emph{Estimated on} column of Table~\ref{tab:monitor} lists. The read-out gain and the named jump come from 25 regime-III units (\#36b--\#51); the field share from 8 regime-III periods; the $\kappa$ table from 9 regime-III periods; and the two relaxation rates of model 16 from \#51 alone. In regime I the read-out gain is near zero (Fig.~\ref{fig:dial}), so the coupling results of model 02 are regime-III results.

% ---------------------------------------------------------------- 02 + 09
\subsection{Model 02: kinetic Ising with read-out coupling (with 09 as gain bookkeeping)}\label{sec:m02}

\textbf{Model.} Eqs.~\eqref{eq:pic-kinetic}--\eqref{eq:fano} of Sec.~\ref{sec:p02}. This form differs from plain kinetic Ising~\cite{aguilera2026} in three ways: the clock is the agent's own calls, the input is the read batch and not the current state of the others, and $J$ depends on whether the message names the reader. The Hawkes view (model 09~\cite{hawkes1971}) adds the gain and Fano bookkeeping.

\textbf{Mapping.} Spins are call outcomes (talk, mention, reply, content move). Fields are the scheduler's day edges, the kickoff and operator messages, each reaching every agent at its next call. Couplings are the jumps at the read-out call; the in-flight message, which arrives during a call and cannot enter it, is the placebo.

\begin{figure}[t]
\includegraphics[width=\columnwidth]{js/h08_readout_jump.pdf}
\caption{The response starts at the read-out call (model 02, H08). (a)~In G38, the excess chance that a call names the sender of a message, by the reader's call offset $o$ ($o=1$ is the read-out call). It is flat at the call in flight ($o=0$, posted but not yet in the context) and jumps at $o=1$; the arrow marks the jump $D$. The other-room placebo (gray) stays at zero. (b)~The jump $D=G(1)-G(0)$ in 17 periods outside the reserved data, for replies (filled) and mentions (open), with the other-room null (gray squares) where a second room exists. Look at the step between $o=0$ and $o=1$ in (a), and at the gray squares at zero in (b).}
\label{fig:readoutjump}
\end{figure}

\begin{figure}[t]
\includegraphics[width=\columnwidth]{js/h67_named.pdf}
\caption{The loop gain is small, and reads that name the reader carry about half of it (model 02, H67). Each bar is a regime-III unit (a goal period, or the part of one between two scaffold changes). The bar height is the read-out loop gain $g$: the number of extra talk calls that one message causes directly. The dark part comes from reads that name the reader, and the light part from all other reads, which are far more numerous. In three units the named part alone exceeds $g$, so there the other reads add nothing. Black lines: 95\% CI (cut at 0). The dashed line at the top is $g=1$, where echoes would run away. Look at the scale: every unit sits far below $g=1$.}
\label{fig:named}
\end{figure}

\textbf{Best-supported cards.}
\begin{enumerate}
\item \textbf{H50: content couples at the read-out call.} In regime III a read moves the reader's content by $J^c_1=0.033$ [0.027, 0.039] (bge; gte 0.046), with the CI above 0 in 10/16 units and below 0 in none. Named messages give 0.077 [0.063, 0.091], unnamed 0.016 [0.009, 0.022], about five times less. It beats the in-flight placebo and a rival without the read-out step, and the content channel was never fitted\cred{0.79}. The scheduler's day edges explain 0.62 (regime I) and 0.67 (regime III) of activity co-movement.
\item \textbf{H111: a sum rule for talk fluctuations.} In regime III the observed Fano ratio of 10-min talk counts matches Eq.~\eqref{eq:fano} with $g$ from H67: $r_F=\Phi_{\rm obs}/\Phi_{\rm pred}=1.01$ [0.94, 1.09] over 18 units, with no between-unit spread (Fig.~\ref{fig:fano}). $r_F$ lies within $\pm20\%$ in 14/18 units, on point estimates. Cross-room covariance is positive in 0/12 units\cred{0.65}. Regime I has an excess, $r_F=1.34$ [1.18, 1.52], whose cause is unknown.
\item \textbf{H99: collective memory is read-out coupling.} Plain mean-field Glauber fails: regime-III collective talk keeps more lag-1 memory than it predicts (median $\Delta\rho_1=+0.052$). The hop-1 gain of H67 reproduces the excess ($g_{1,\rm fit}/g_{\rm lag}=0.96$, 24 units)\cred{0.63}.
\item \textbf{H08: the response starts at the read-out call.} The chance that a recipient names or answers the sender jumps at the read-out call and not at the call in flight (Fig.~\ref{fig:readoutjump}): 14/17 periods for mentions, 17/17 for replies, with the other-room placebo flat in 8/8 two-room periods (Fig.~\ref{fig:readout}). Newcomers name an old-timer only after they receive its message (35/35 pairs). Content moves only inside replies that name or answer the sender: $+1.25$ [$+0.98$, $+1.60$] cos$\times100$ with them, $-0.39$ [$-0.62$, $-0.12$] without (G51)\cred{0.85}. Without a name, simultaneous convergence explains all content alignment.
\item \textbf{H40: the clock is the call.} In regimes II and III the reply hazard given a talk call has elasticity $\eta_{\rm rep|talk}=-0.11$ [$-0.16$, $-0.06$] to call duration, and the wall clock ($\eta=1$) is excluded in every period. The call clock beats the wall clock on held-out data in 33/33 periods\cred{0.82}. Regime I does not follow: $\eta=+0.51$ [0.41, 0.61], unexplained.
\item \textbf{H42: a delayed step at the next read-out.} A read that names the recipient adds 0.076 [0.066, 0.086] talk at its read-out call (26 units); half of it, 0.046 [0.037, 0.056], comes with no exchange in progress, so the name itself acts. The named kernel keeps 0.90 of its size against a fitted shared (Cox) field, while a generic exponential cross term keeps 0.22\cred{0.80}. With per-agent call clocks as the baseline, wall-time cross-excitation falls from 0.061 to $n_{\rm cross}=0.004$.
\end{enumerate}

\begin{figure}[t]
\includegraphics[width=\columnwidth]{js/h111_fano.pdf}
\caption{A sum rule that links fluctuations to response for talk (models 02 and 09, H111; Eq.~\eqref{eq:fano}). (a)~In a simulated linear Hawkes swarm, the collective Fano ratio $\Phi$ of 10-min talk counts follows $1/(1-g)^2$ (room-adjusted for 15 agents) with no free parameter. (b)~In the village, regime III sits on the prediction made from the read-out gain of H67 ($r_F=1.01$ [0.94, 1.09], 18 units). Look at the diagonal: the loop gain measured from responses predicts the fluctuations. Regime I lies above it ($r_F=1.34$ [1.18, 1.52]), for a reason we cannot explain.}
\label{fig:fano}
\end{figure}

\textbf{Variant: the context-held self-field.} An agent's own recent output, held in its context window, acts as a field on itself; inside idle traps, partly like a P\'olya urn. This variant replaces Kramers escape (H16) and carries four cards. For project choice the context acts as a reset step, not as a proportional urn: the own-context share does not set the odds of leaving a project (urn slope $+0.07$ [$-0.05$, $+0.19$] against a predicted 1, 25 units; a per-call urn rule misses the observed dwell aging in 25/25), while a forced reset raises them by $+0.40$ [$+0.25$, $+0.54$] log-odds in G51, against a placebo of $-0.02$ (H134). Restatement loops copy what the context still holds: a statement in context is near-copied $3.38\times$ [2.56, 4.46] more often than an erased one at equal lag (4/4 periods), and an erasure ends loops as a step, not in proportion to the text removed (H69)\cred{0.88}. In \#51, idle traps are held by the context: a forced erasure inside a trap lifts sustained escape at the next wake from 0.47 to 0.84 (log OR $+2.68$ [2.31, 3.16]), and a directed read frees an agent by address, not by dilution ($+0.45$ [0.32, 0.58]) (H16)\cred{0.68}. The context's own idle share carries 65\% [59, 71] of the held-out information of the aging clocks; input starvation carries none (H72)\cred{0.65}. A forced erasure gives a one-call re-reading spike ($+0.106$ [0.089, 0.122], 9/9 periods) with a tail of 8.2 [5.8, 11.8] calls and a write dip of 20--32\%, by reference window (H44)\cred{0.85}.

\textbf{What failed.} Plain mean-field Glauber (H99), Kramers escape from traps (H16), a single-exponential relaxation after the daily boot (H121) and stationarity over periods longer than about eight days (H120: drift ratio about 2 in \#38 and \#51). Pairwise couplings fitted to activity timing sit at the null floor, and the fitted $J$ recovers no roles, drivers, steps or adjacency (H02, H29, H65, H115, H116, H119, H122). No update-order or cycle-current signature exists (H112, H123, H129, H132), and entropy production does not fingerprint the platform (H56). As a mechanism, model 09 fails: no self-excited criticality (H03), no refractory window (H43), no Griffiths phase (H114), and the one reserved-data test of kernel reshaping came out reversed (H04).

\textbf{Verdict: useful}, in regime III and in its read-out form. It is the positive core of the paper. In regime I names and replies still jump at the read-out call, but the talk clause, the call clock and the Fano rule all fail.

% ---------------------------------------------------------------- 11 + 01
\subsection{Model 11: vector spins as a field model (with the 01 mean-field gauge)}\label{sec:m11}

\textbf{Model.} Eq.~\eqref{eq:vector} of Sec.~\ref{sec:p11}. The inverse Ising model (01) is the binary, mean-field gauge of the same picture: with magnetization $m=\tanh[\beta(J_0m+h)]$, the susceptibility is $\chi\propto1/(1-g)$ with $g=\beta J_0(1-m^2)$~\cite{schneidman2006,meshulam2019}. Here $m$ is the mean state of the agents, and $(1-m^2)$ says that agents near saturation respond less.

\textbf{Mapping.} The goal is a literal field vector: the embedding of the kickoff text. Rooms, style and private roles are static fields. Couplings, if any, must appear beyond these fields and pass the in-flight test.

\begin{figure}[t]
\includegraphics[width=\columnwidth]{js/h54_kickoff_rank.pdf}
\caption{The own kickoff is the target (model 11, H54). (a)~Rank of the period's own kickoff text among 33 candidates, by the similarity of Fig.~\ref{fig:kickoff}: first in 18/33 periods, against a swap null whose median rank is 17 (gray band: 90\%). Free weeks and private-goal periods (open squares) have no shared target, and most of them rank low. (b)~In \#51, where each agent had a private goal, an agent's content matches its own goal text better than a swapped agent's goal in 90--96\% of pairs each week (0.95 over all weeks; chance 0.5). Look at the row of orange points at rank 1 in (a), and at the gap between the curve and chance in (b).}
\label{fig:kickoffrank}
\end{figure}

\textbf{Best-supported cards (11).}
\begin{enumerate}
\item \textbf{H54: the kickoff text is the target.} Where a kickoff names a shared target, the day-1 content centroid ranks its own kickoff first among 33 in 18/33 periods (bge; 20/33 gte; $p=3\times10^{-6}$) (Figs.~\ref{fig:kickoff} and \ref{fig:kickoffrank}). In \#51 each agent sits on its own private goal (0.95). A human message pulls each reader by $\Delta\approx0.09$ at its first message after the read, and reading adds $+0.11$--0.12 over unread in-flight messages\cred{0.90}.
\item \textbf{H96: a goal switch is a quench, not a hysteresis loop.} On day 1 the old state's field-orthogonal persistence is 0.18--0.27, against 0.85--0.89 across an ordinary night ($\Delta R_1=-0.62$ [$-0.78$, $-0.45$], bge), below the night reference in 20--21 of 23--24 transitions\cred{0.74}. The coercivity test has power 0.50, so hysteresis is untested rather than excluded.
\item \textbf{H13: family fields are mostly style, and labs couple a little more.} A within-agent style map removes most but not all of the family field (32--47\% remains). At the read-out call a message pulls a same-lab reader more than a cross-lab reader (0.067 vs 0.039; $\Delta=0.029$ [0.010, 0.054], 8/8 regime-III units); about half of $\Delta$ is lab susceptibility. Talk has no family term\cred{0.74}.
\item \textbf{H46: style is an identity charge with a context-held offset.} Style plus function words identify agents across goal switches at 0.85 (content 0.51, chance 0.15). A forced erasure moves style (agent-scaled percentile 0.546 [0.522, 0.564]) while content stays put (0.51), because each erasure redraws a small per-segment offset\cred{0.79}.
\item \textbf{H100: rooms split content spontaneously} (Sec.~\ref{sec:cross}).
\item \textbf{H82: the previous goal persists in members.} Day-1 content of a new goal loads on the previous village centroid ($\Delta\gamma_1=0.158$ against a null 95th percentile of 0.055, 25 boundaries); each agent's own previous content absorbs the trace (post hoc)\cred{0.73}.
\end{enumerate}

\textbf{Best-supported cards (01 gauge).} H67: the read-out talk loop gain is subcritical in 66/66 units, regime-III median $g=0.13$ (Fig.~\ref{fig:named}; pool 0.127, 25 units; units 0.00--0.29; every upper bound below 0.40). It agrees with the gain implied by fluctuations, $g_{\rm Fano}=0.149$ [0.123, 0.175]; this is a consistency check, not an independent test. On the chat clock regime I has no read-out gain ($-0.022$ [$-0.052$, 0.009]), and an earlier gain twice as large (H50) was an estimator artifact\cred{0.87}. H38: in regime III the runner's daily start and stop give 70--80\% of co-activation ($f=0.72$ median over 8 periods; 0.83 after trimming; regime I 0.11). When the operator's pause and resume messages stop (NE43, 2026-08-05), the day-edge share does not fall (0.189 to 0.135, $p=0.08$), so the edges are the runner's schedule\cred{0.90; fragile}. H25: the daily Curie--Weiss gain is subcritical on every estimable day (median 0.16 for activity and talk)\cred{0.85}. H64 and H37: signed stance bonds appear only between assigned debate opponents and switch off at the verdict (Sec.~\ref{sec:cross})\cred{0.64, 0.65}.

\textbf{What failed.} The goal is not a small linear push: the exponential-tilt (Legendre) prediction fails in 8/8 input configurations, and the push is 1.4--3.1 free-week standard deviations (H10). Content is not near-critical (H26). Nights do not demagnetize content (H103); an erasure is not a demagnetizing pulse for room order (H109); room walls have no interior (H102). Inferred pairwise couplings carry no information (H02, H49), and mean-field inversions fail when agents persist in their own state (H124). Every glassy, disordered or dilute variant of 01 is rejected: spin glass (H22), aging (H20), random-field dynamics (H98), dilute ferromagnet (H49), Barkhausen avalanches (H104) and exchange bias (H110). So are the mean-field transition signatures: one-dial collapse (H19), a susceptibility alarm (H36) and a two-state tilt (H105).

\textbf{Verdict: useful as a field model}; rejected as a coupled ordered magnet. Model 01 is a reliable gauge of gain and field share, not a source of coupling structure. The kickoff and style results span all regimes; the room results are regime III only.

% ---------------------------------------------------------------- 16
\subsection{Model 16: Langevin relaxation}\label{sec:m16}

\textbf{Model.} Eq.~\eqref{eq:pic-langevin} of Sec.~\ref{sec:p16}, on the call clock. An underdamped well, $\ddot x+2\zeta\omega_0\dot x+\omega_0^2(x-c)=\xi$ with damping ratio $\zeta<1$, would overshoot and then undershoot after a step; H125 tests this rival.

\textbf{Mapping.} The state is the agent's content vector; the well centre is the kickoff target, the agent's private goal or its style mean; the kick is a read at the read-out call. Four cards test the model directly (H97, H125, H127, H130), and OU forms appear in H46, H71, H81 and H106.

\textbf{Best-supported cards} (Fig.~\ref{fig:m16}, Sec.~\ref{sec:p16}).
\begin{enumerate}
\item \textbf{H97: a kickoff is an anisotropic restoring force.} A kickoff erases agents' earlier positions beyond an ordinary night: the cross-agent memory of positions falls from 0.68 to 0.42, an extra forgetting $\Delta\rho=+0.26\pm0.04$, positive in 16/18 kickoffs and in both embedding models. The well is stiffer along the kickoff direction (0.63) than across it (0.25), and day 1 overshoots ($+0.12\pm0.02$)\cred{0.70}.
\item \textbf{H125: the response is overdamped.} After a kickoff, alignment overshoots on day 1 ($E_1=+0.055\pm0.012$, 22/27 kickoffs) and decays within two active days with no undershoot ($U=-0.006$ [$-0.020$, $+0.008$], 90\% CI; power 1.00 at $\zeta\le0.5$). The pre-registered oscillator fails, and the overshoot is a field that fades, not inertia\cred{0.64}. Damping between 0.5 and 1 is not excluded.
\item \textbf{H130: a fast read kick on a slow well.} In \#51 a read pulls the reader's next statement toward the sender beyond the in-flight placebo ($J_K=0.044$ [0.038, 0.050], 11/12 units). The pull decays in about 7 calls ($\gamma_k\approx0.13$--0.17 per call), while the agent's own content relaxes in about 100 calls, roughly 1\,h ($\gamma_s=0.0094$ [0.0076, 0.0115] per call). The rate ratio is 14.8 [7.2, 30.4] (90\%)\cred{0.64}. The pre-registered one-rate model failed its kill rule; the two-rate form was found after the test.
\item \textbf{H81 and H71: slower layers.} Regime-I content carries a collective slow mode with $\tau\approx23$--28 days (Sec.~\ref{sec:m17})\cred{0.61}. Agent memory size returns to an agent-specific set point ($\phi=0.67$ [0.63, 0.71] per consolidation cycle, 37/37 units) with no overshoot, but with two timescales, not one (post hoc)\cred{0.74}.
\end{enumerate}
H127 supports the same picture from the other side: 74\% of 152 rising agents reach their day-1 alignment within the read-out call of the kickoff, and the excess then decays from $2.4\times$ to $1\times$ the plateau over about 64 calls. There is no per-agent catch-up time for the call rate to set\cred{0.59}.

\textbf{What failed.} Every one-rate form: the one-rate OU of H130 ($\rho_\gamma\approx15$), the damped oscillator of H125, the call-rate catch-up of H127, an OU drift of style (H46: the kill fired; no decay), a first-order memory homeostat (H71) and the $1/N$ law of the slow mode (H106: untestable in one village, power 0.00).

\textbf{Verdict: partly useful.} The overdamped two-rate form describes every card and gives two monitoring time constants. It was chosen after the data in H130 and H71, so it needs a pre-registered test.

% ---------------------------------------------------------------- 17
\subsection{Model 17: collective modes and random-matrix spectra}\label{sec:m17}

\textbf{Model.} Eq.~\eqref{eq:spike} of Sec.~\ref{sec:p17}. Dynamic modes come from a VAR(1) or OU fit, $x(t+1)=Ax(t)+\varepsilon$, with lifetimes $\tau_k=-\Delta t/\ln|\mu_k|$ for the eigenvalues $\mu_k$ of $A$; each is one process of model 16.

\textbf{Mapping.} Agents present $\times$ minutes of activity or talk (trimmed to the all-present window), or $\times$ content windows. Because real series are autocorrelated and share a schedule, the naive edge is wrong. We use a calibrated edge: the 95th percentile of $\lambda_1$ under surrogates that keep each agent's own series but break alignment between agents~\cite{laloux1999}.

\textbf{Best-supported cards} (Fig.~\ref{fig:m17}, Sec.~\ref{sec:p17}).
\begin{enumerate}
\item \textbf{H12: talk and content modes are real; the activity mode is the scheduler.} Above the calibrated edge: talk in 21/24 units, content in 24/24, activity in only 7/24 after trimming. A \#12 debate motion compresses content (PR $-25\%$, 7/9 debates, $p=0.010$); kickoffs do not. The lowest-PR weeks are restatement loops\cred{0.77; fragile}. The round-1 activity mode was a data artifact.
\item \textbf{H92: cleaning helps forecasts but ties shrinkage.} Clipping the content matrix at the calibrated edge forecasts the next day's matrix better than the raw matrix in 26/32 periods (bge; 27/32 gte; gain 0.10 [0.03, 0.16]), and the gain grows with $q$. It ties constant-correlation shrinkage within 1--2\%, so the persistent structure is one uniform mode\cred{0.71}.
\item \textbf{H81: one slow mode, in regime I only.} Residual content of blocks up to 14 days apart aligns beyond a composition null and decays with $\tau=28\pm5$ days (bge; $23\pm4$ gte); the mode survives member turnover and does not follow the artifact record\cred{0.61}. Unmeasured calendar drift remains a rival, and no such mode exists in \#51.
\end{enumerate}

\textbf{What failed.} Eigenvector rotation is a poor reorganization alarm (AUC 0.56 [0.40, 0.71] against 0.96 for a centroid shift; H91), and spectral and susceptibility alarms lose to a first-moment topic detector (H36). Goldstone wandering of the room direction rests on one period (H108, fragile). The spiked-mode predictions of Eq.~\eqref{eq:spike} ($\lambda_1$ at other window lengths, split-half overlaps) and dynamic mode lifetimes are untested.

\textbf{Verdict: partly useful}, as a gauge of how many collective modes exist and how uniform they are. Its sharpest predictions are untested.

